\subsection{Four and five rows with five outside elements}

\begin{lemma}[A missing old column character]
\label{C:missing-column-character-act-cups}
Let $M$ have a minimal ideal with two columns and row
augmentation module $V$.  If
$\operatorname{Hom}_A(V,\mathbb Q_\chi)=0$ for its old
column character $\chi$, all rational degree-two products
of every finite act cone are zero.
\end{lemma}
\begin{proof}
After the constant-action quotient and compression, each
balanced form $L$ restricts on the old column augmentation
line to a covector $\ell_L$ satisfying
$\ell_L A_m=\chi(m)\ell_L$.
The hypothesis makes it zero.  Consequently the honest
transferred affine cocycle has
$b^L_m=\beta(m)\ell_L=0$ on every old actor, and its value
on every pair of old actors is zero.  The second pair in
any nondegenerate four-chain of the act category consists
of old actors.  Thus every degree-two product is already
zero as a cochain on the quotient cone.  Pullback along
the degree-two surjection proves the result.
\end{proof}

\begin{theorem}[Four rows and five outside elements]
\label{C:four-row-five-outside-seven-states}
Let $M$ have minimal band $4\times2$ and exactly five outside
elements.  No finite right act with at most seven states
gives a cone of type $\mathbb CP^2$.
Equivalently, this entire minimal-band/outside-size branch
is excluded through twenty-one category arrows.
\end{theorem}
\begin{proof}
Assume such a cone exists.  A state fixed by the minimal band
would make its degree-two products zero.  Otherwise use
Lemmas~\ref{C:two-column-projective-orbit-kernel} and
\ref{C:two-column-finite-quotient-model}.
If $o$ is the number of two-point core orbits, the compressed
monoid $N$ has a finite four-dimensional model and band--Tor
modules with
\[
 \dim V=3,\quad \dim W=|S|-2o+1,\qquad
 \bigl(\dim\operatorname{Tor}_0,
       \dim\operatorname{Tor}_1,
       \dim\operatorname{Tor}_2\bigr)=(o,0,1),
\]
and zero higher Tor, over the same five-dimensional contracted
algebra $A$ as $M$.
If $A$ is local, Lemma~\ref{C:local-two-column-act-cups}
already gives a contradiction.  Otherwise
Proposition~\ref{B:five-nonlocal-exceptions} leaves the line
and Nakayama algebras of Lemma~\ref{B:five-exceptional-targets}.

For the Nakayama algebra, its exact indecomposable Tor table
shows that the only three-dimensional left module compatible
with $\operatorname{Tor}_1=0$ and one-dimensional
$\operatorname{Tor}_2$ is
\[
 V=S_2\oplus Ae_2.
\]
Indeed a degree-two contributor requires exactly one copy of
$S_2$ or $L$; any added $S_1$ forces degree-one Tor, and
$L\oplus S_2$ contributes twice.  The remaining projective
and semisimple cases contribute zero or an even multiplicity.
The only one-dimensional quotient character of the displayed
module is $S_2$, and it occurs twice.
If the old column character differs, use
Lemma~\ref{C:missing-column-character-act-cups}; if it agrees,
use Lemma~\ref{C:two-covector-affine-origin}.
Both contradict the projective-plane square.

For the line algebra write
$L_a=Ae_2$, $L_b=Ae_3$, $R_a=e_1A$, $R_b=e_2A$.
The degree-two condition requires one $S_3$ in $V$ and one
$T_1$ in $W$; the degree-one condition excludes $S_2$ from
$V$ and $T_2$ from $W$.
Thus the only possible row modules are
\[
 2S_1\oplus S_3,\qquad S_3\oplus L_a,
 \qquad S_3\oplus L_b.
\]
The first is split into characters and is excluded by
Proposition~\ref{C:split-row-act-cups}.
The last has only the quotient character $S_3$, with
multiplicity two, so the preceding affine argument excludes it.

For $V=S_3\oplus L_a$, write
\[
 W=T_1\oplus T_3^{\oplus s}
       \oplus R_a^{\oplus p}\oplus R_b^{\oplus q}.
\]
The exact tensor dimensions give
\[
 o=s+p+q,\qquad
 1+s+2p+2q=|S|-2o+1,
 \qquad |S|=4o-s.
\]
Since $s\leq o$ and $|S|\leq7$, one has $o\leq2$.
For $o=1$, $|S|\leq4$, already excluded by the small-state
theorems.  For $o=2$ the only cases are
\[
 \begin{array}{c|c|c}
 |S|&s&(p,q)\\ \hline
 6&2&(0,0)\\
 7&1&(1,0)\text{ or }(0,1).
 \end{array}
\]
For six states, both degree-two forms factor through the
single $S_3$ row summand, so they have matrix rank one.
For seven states with $(p,q)=(0,1)$, the contributions from
$T_3\otimes S_3$ and $R_b\otimes L_a$ give a basis of two
rank-one balanced forms.  Every basis class has square zero
by Lemma~\ref{B:rank-one-square}.
Their pullbacks span the one-dimensional $H^2$ of the original
cone.  At least one has a nonzero pullback, contradicting the
projective-plane square.  This uses a spanning set of zero
squares, not an unsupported claim about their mixed product.

For the remaining pair
\[
 V=S_3\oplus L_a,
 \qquad W=T_1\oplus T_3\oplus R_a,
\]
both representation kernels in $A$ are exactly $\mathbb Qb$.
The complete actual-basis result
Lemma~\ref{B:five-line-bases}, applied to the opposite algebra
if necessary, writes $b$ as a nonzero multiple of either one
distinguished basis element or a difference of two such elements.
In the first case both augmentation operators are zero, so
both finite set maps are constant and the actor profile lies
in the band.  In the difference case the two augmentation
operators agree on each side and are nonzero: otherwise two
independent basis elements would lie in the one-dimensional
kernel $\mathbb Qb$.  Equality of nonzero augmentation
operators determines equality of the actual finite set maps.
Thus the faithful profile quotient of $N$ kills $b$ modulo
the band.  Its contracted algebra is a quotient of
\[
 A/(b)\cong\mathbb Q\times\mathbb Q(1\longrightarrow2).
\]
Every quotient of this particular one-arrow algebra is
hereditary.  Its fourth rational homology is zero, whereas
the natural map preserves the balanced-tensor description
of $H^2$.  All degree-two products on $N$ pull back from
zero products in that quotient, a final contradiction.
\end{proof}

\begin{lemma}[The four-by-four module dimensions at five outside elements]
\label{C:five-by-five-five-outside-CP2}
A monoid with minimal band $5\times5$ and exactly five
outside elements cannot have type $\mathbb CP^2$.
\end{lemma}
\begin{proof}
The row and column augmentation modules both have dimension
four and require Tor dimensions $(1,0,1,0,\ldots)$ over a
five-dimensional algebra.  The local case is excluded by
Proposition~\ref{B:local-quotient-CP2}.  The nonlocal
classification leaves the same two algebras as above.

For the line algebra the degree-one and degree-two conditions
give exactly the following row and column possibilities.
Their degree-zero Tor dimensions are
\[
\begin{array}{c|ccc}
 &T_1\oplus3T_3&T_1\oplus T_3\oplus R_a&
                  T_1\oplus T_3\oplus R_b\\ \hline
 3S_1\oplus S_3&6&7&4\\
 S_1\oplus S_3\oplus L_a&4&4&3\\
 S_1\oplus S_3\oplus L_b&7&4&4
\end{array}
\]
None is one.
For the Nakayama algebra the only remaining left modules are
$S_2\oplus Ae_1$ and $L\oplus Ae_2$.
For the former, the admissible four-dimensional right modules
with one degree-two contribution are
$T_2\oplus e_1A$ and $R\oplus e_2A$;
for the latter, only $T_2\oplus e_1A$ remains.
Their tensor dimensions are all three, not one.

These lists can also be checked directly from the full
indecomposable table: \path{code/a_five_four_by_four_modules.py}
enumerates all thirty dimension-four modules on each side
for the line and all sixteen on each side for the Nakayama
algebra, checking all $900+256$ pairs.
The exact output is
\path{code/outputs/a_five_four_by_four_modules.json}; it has
no pair with Tor triple $(1,0,1)$.
\end{proof}

\begin{corollary}[Five rows and five outside elements]
\label{C:five-row-five-outside-five-states}
Let $M$ have minimal band $5\times2$ and five outside elements.
No act with at most five states gives a projective-plane cone.
Thus this entire branch is excluded through twenty-one arrows.
\end{corollary}
\begin{proof}
At most three states are already excluded below twenty-two
arrows.  For four states the band dichotomy gives zero products
or constant action.  In the latter case compression has band
$5\times4$ with five outside elements, excluded by
Lemma~\ref{C:four-by-five-five-outside-CP2}.
For five states, a fixed state again gives zero products.
Otherwise there are one or two two-point core orbits.
With one, the action is constant and compression has band
$5\times5$, excluded by the preceding lemma.
With two, the projective-kernel and finite-model lemmas give
a compressed monoid with band $5\times3$, degree-two rank
two and fourth homology $\mathbb Z$.
The rank-two part of
Theorem~\ref{B:three-by-five-targets}, after opposition,
makes every rational degree-two product zero on this finite
source.  Its degree-two classes pull back onto those of the
original cone, contradicting its nonzero square.
\end{proof}
